In Twenty20 (T20) cricket, the death overs --- the final five overs (16--20) of
an innings --- carry a disproportionate share of the outcome. Scoring rates peak,
wickets are spent in pursuit of boundaries, and the required run rate in a chase
swings sharply on a single delivery. Two intuitions dominate how commentators,
coaches, and models reason about this phase. The first is that \emph{who} is on
strike or bowling --- a death-overs finisher against a yorker specialist ---
largely determines the next ball. The second is that \emph{momentum} matters:
that a run of dot balls builds pressure that spills into the next delivery, or
that a batter who has just found the boundary is ``hot.'' Both intuitions are
widely held; neither has been tested cleanly at the level of an individual
death-over delivery.

The quantitative cricket literature provides the machinery to test them but has
not done so for this regime. Per-ball win-probability and expected-runs models
--- WASP's dynamic-programming value function \citep{brooker2012wasp}, the
dynamic logistic regression of \citet{asif2016inplay}, and the hierarchical
empirical-Bayes simulators of \citet{swartz2016simulator, swartz2020xer} ---
either deliberately exclude player identity (``average team vs.\ average team'')
or include it for simulation rather than for a clean state-versus-identity
attribution. The pressure-index family \citep{shah2014pressure,
bhattacharjee2016quantifying, mallawa2024pressure} \emph{defines} pressure as a
deterministic function of game state and then validates it by predicting
outcomes; it never isolates a residual, history-beyond-state component. And the
near-universal Markov assumption of these models --- that the next ball depends
only on the current state, not on the path that produced it --- has not been
tested against pressure sequences at the delivery level. Meanwhile, the
general-sports ``hot hand'' literature has been transformed by the discovery of
a streak-selection bias in the naive conditional-probability estimator
\citep{miller2018surprised, miller2024coldshower}, a correction that overturned
the canonical basketball null \citep{gilovich1985hothand}. No cricket study has
imported it.

We address all three gaps in a single experiment on real Cricsheet data. We
build calibrated per-ball models of expected runs $E[\text{runs}]$ and wicket
probability $P(\text{wicket})$ for the T20 death overs, decompose their
predictive skill into game state and batter/bowler identity, test whether
pressure history adds skill beyond a Markov state model, and apply the
Miller--Sanjurjo correction to the question of per-ball momentum.

We find, first, that the apparent death-over ``cold hand'' is largely an artifact
of the estimator rather than the data. By importing the Miller--Sanjurjo correction
--- new to cricket --- we measure how large that artifact is: the naive momentum
statistic ($D=-0.114$) is almost exactly cancelled by the finite-sequence bias
($-0.121$), so we put the size of the apparent-momentum illusion at $+0.121$ and
the corrected, state-controlled estimate at a non-significant $+0.007$. Quantifying
the illusion, not merely reporting a null, is our central and most defensible
contribution. Second,
we have no evidence that a \emph{historical-rate encoding} of batter/bowler
identity improves held-out per-ball skill over a game-state model on either target
--- and, more pointedly, that adding this identity encoding \emph{hurts}
out-of-sample expected-runs accuracy in every one of six rolling held-out seasons,
even as SHAP attributes the majority of explanatory mass to batter identity. A
random-player placebo confirms this is a genuine dissociation between faithful
in-sample attribution and generalizable skill, not a leakage or feature-encoding
artifact. Third, residual dot-ball pressure adds significant skill to expected runs
(with the expected negative sign) but nothing to wicket risk, so the Markov
property is violated only for scoring intensity, not for dismissal risk. The
absolute proper-score differences here are small --- per-ball death-over outcomes
sit close to the base rate, so even the best models improve on it only modestly ---
and we therefore lean on the qualitative dissociation and the bias correction
rather than the size of any one model gain. All conclusions are specific to per-ball
outcomes in the death overs, under our feature set and our historical-rate identity
encoding, after state control and bias correction; we are careful throughout to
distinguish ``this encoding of identity does not help'' from ``identity does not
matter,'' and ``no significant momentum'' from ``momentum proven absent.''

\paragraph{Contributions.} We make the following contributions:
\begin{itemize}
  \item To our knowledge, we provide the first application of the Miller--Sanjurjo
        streak-selection correction to cricket, and use it to \emph{quantify} the
        apparent-momentum illusion at the death over: the $+0.121$ gap between the
        naive momentum estimate ($D=-0.114$) and the finite-sequence bias
        ($-0.121$). The corrected estimate ($+0.007$) is not significant against a
        state-stratified null ($p=0.40$), so the apparent ``cold hand'' is an
        estimator artifact, and we report how large that artifact is.
  \item We document a SHAP-versus-skill dissociation and stress-test it: TreeSHAP
        credits batter identity with $51$--$64\%$ of explanatory mass, yet our
        historical-rate encoding of that identity adds zero or negative held-out
        incremental skill (runs RMSE $-0.0053$, 95\%
        CI $[-0.0089, -0.0019]$). A random-player placebo collapses batter SHAP mass
        to $7$--$8\%$ while leaving skill unchanged, and a six-season rolling-window
        evaluation finds identity significantly hurts expected-runs RMSE in
        \emph{every} season --- showing the gap is genuine overfitting of
        non-stationary player form, not contamination or a single-season fluke.
  \item We test the Markov assumption directly and find it violated for
        $E[\text{runs}]$ only: pressure history adds significant skill (RMSE
        $+0.0065$, 95\% CI $[0.0048, 0.0083]$, dot-run sign negative) but nothing
        to $P(\text{wicket})$.
  \item We release a calibrated per-ball benchmark for T20 death overs on
        223{,}678 deliveries from 4{,}402 men's IPL and T20I matches
        (2005--2026, held-out 2026 test season), reporting proper scores for
        both $E[\text{runs}]$ (RMSE $1.728 \to 1.686$ over the base rate) and
        $P(\text{wicket})$ (log-loss $0.2748 \to 0.2738$). Because per-ball
        death-over outcomes sit close to the base rate, even the best model improves
        on it only modestly; we treat this low ceiling on per-ball predictability as
        a substantive finding and therefore rest our claims on qualitative
        dissociations and the bias correction rather than on the magnitude of any
        single proper-score gain.
\end{itemize}

The remainder of the paper reviews related work (Section~\ref{sec:related}),
formalizes the targets and methods (Sections~\ref{sec:background}
and~\ref{sec:method}), reports results (Section~\ref{sec:experiments}), and
discusses interpretation, limitations, and future work
(Section~\ref{sec:discussion}).
